\documentclass{article}
\usepackage[utf8]{inputenc}
\usepackage{amsmath,amssymb}
\usepackage[most]{tcolorbox}
\usepackage{geometry}
\geometry{margin=2.5cm}

\newcommand{\step}[1]{\par\noindent\hangindent=3.5em\hangafter=1%
  \makebox[3.5em][l]{\textbf{#1.}}}
\newcommand{\steptag}[1]{\hfill{\small\textsf{[#1]}}}

\newtcolorbox{proofbox}[1]{
  colback=gray!5, colframe=gray!60,
  fonttitle=\bfseries, title={#1},
  breakable, enhanced,
  left=4pt, right=4pt, top=4pt, bottom=4pt,
  before skip=10pt, after skip=10pt
}

\begin{document}

\begin{proofbox}{For r in A the square hole q\_r = P(r\textasciicircum{}2)-r\textasciicircum{}2 in Ker P sati...}
\step{1} For r in A the square hole q\_r = P(r\textasciicircum{}2)-r\textasciicircum{}2 in Ker P satisfies q\_r >= -C eta ||r||\textasciicircum{}2 1, ||P(q\_r\textasciicircum{}2)|| <= C eta ||r||\textasciicircum{}4, and ||q\_r||\_omega <= C sqrt(eta) ||r||\textasciicircum{}2. \steptag{validated} \\[6pt]
\step{1.1} q\_r in Ker P. Let r in A=Im P (def-near-fixed-algebra), and let q\_r=P(r\textasciicircum{}2)-r\textasciicircum{}2 be its square hole (def-square-hole), where r\textasciicircum{}2=r o r is the ambient Jordan square in V=B(H)\_sa. By GT-Pprops, P is real-linear and P\textasciicircum{}2=P. Since r in Im P, r=P(s) for some s, so P(r)=P\textasciicircum{}2(s)=P(s)=r, i.e. P(r)=r. Then by linearity P(q\_r)=P(P(r\textasciicircum{}2)-r\textasciicircum{}2)=P\textasciicircum{}2(r\textasciicircum{}2)-P(r\textasciicircum{}2)=P(r\textasciicircum{}2)-P(r\textasciicircum{}2)=0 (using P\textasciicircum{}2=P on the first term). Hence q\_r in Ker P. \steptag{validated} \\[4pt]
\step{1.2} LOWER BOUND: q\_r >= -C eta ||r||\textasciicircum{}2 1. Write delta=||P-Phi|| (GT-Pprops), with delta<=C eta. STEP A (r is Phi-near-fixed): since r in A=Im P, P(r)=r (node 1.1), so ||Phi(r)-r||=||Phi(r)-P(r)||=||(Phi-P)(r)||<=delta||r|| (GT-Pprops, delta=||P-Phi||<=C eta). STEP B (Jordan-Schwarz): Phi is unital positive on B(H)\_sa, so by GT-jordan-schwarz with self-adjoint a=r (so a o a=r\textasciicircum{}2 and (Phi r) o (Phi r)=Phi(r)\textasciicircum{}2), Phi(r)\textasciicircum{}2 <= Phi(r\textasciicircum{}2), i.e. Phi(r\textasciicircum{}2) >= Phi(r)\textasciicircum{}2. STEP C (perturb the square): put x:=Phi(r) in V=B(H)\_sa. The Jordan product is commutative, so x o x - r o r = (x-r) o (x+r) (expand: (x-r) o (x+r)=x o x + x o r - r o x - r o r = x\textasciicircum{}2 - r\textasciicircum{}2, using commutativity x o r=r o x). Bounds: ||x-r||=||Phi(r)-r||<=delta||r|| (Step A). For ||x+r|| we use only the in-scope norm triangle inequality and the grounded delta bound (NO ||Phi||=1): since r=P(r) in A, ||Phi(r)-r||=||(Phi-P)(r)||<=delta||r|| (Step A, GT-Pprops), so by the norm triangle inequality ||Phi(r)||=||(Phi(r)-r)+r||<=||Phi(r)-r||+||r||<=delta||r||+||r||=(1+delta)||r||; hence max(||Phi(r)||,||r||)<=(1+delta)||r|| and ||x+r||<=||x||+||r||=||Phi(r)||+||r||<=(1+delta)||r||+||r||=(2+delta)||r||. By GT-jb-submult (||a o b||<=||a|| ||b|| for the Jordan product), ||Phi(r)\textasciicircum{}2-r\textasciicircum{}2||=||(x-r) o (x+r)||<=||x-r|| ||x+r||<=delta||r||*(2+delta)||r||=(2+delta) delta||r||\textasciicircum{}2=:C delta||r||\textasciicircum{}2 (C=2+delta<=2+delta\_0 universal \& dimension-free, delta<=delta\_0=C eta\_0). By GT-orderunit-def, ||z||<=c (z=z\textasciicircum{}*) gives -c1<=z<=c1, so with z=Phi(r)\textasciicircum{}2-r\textasciicircum{}2: Phi(r)\textasciicircum{}2 >= r\textasciicircum{}2 - C delta||r||\textasciicircum{}2 1. Combining with Step B: Phi(r\textasciicircum{}2) >= Phi(r)\textasciicircum{}2 >= r\textasciicircum{}2 - C delta||r||\textasciicircum{}2 1. STEP D (replace Phi by P): ||(P-Phi)(r\textasciicircum{}2)||<=delta||r\textasciicircum{}2||=delta||r||\textasciicircum{}2 (GT-bh-cstar ||r\textasciicircum{}2||=||r||\textasciicircum{}2 for r=r\textasciicircum{}*), so by GT-orderunit-def P(r\textasciicircum{}2) >= Phi(r\textasciicircum{}2) - delta||r||\textasciicircum{}2 1 >= r\textasciicircum{}2 - (C delta+delta)||r||\textasciicircum{}2 1 = r\textasciicircum{}2 - C delta||r||\textasciicircum{}2 1 (absorb constants). Hence q\_r=P(r\textasciicircum{}2)-r\textasciicircum{}2 >= -C delta||r||\textasciicircum{}2 1 >= -C eta||r||\textasciicircum{}2 1 (delta<=C eta). Uses node 1.1. \steptag{validated} \\[4pt]
\step{1.3} BOUND ||Phi(q\_r\textasciicircum{}2)|| and ||P(q\_r\textasciicircum{}2)|| <= C eta ||r||\textasciicircum{}4. Write delta=||P-Phi||<=C eta (GT-Pprops). STEP A (crude norm bound on q\_r): q\_r=P(r\textasciicircum{}2)-r\textasciicircum{}2, so ||q\_r||<=||P(r\textasciicircum{}2)||+||r\textasciicircum{}2||<=||P|| ||r\textasciicircum{}2||+||r||\textasciicircum{}2<=(1+C eta)||r||\textasciicircum{}2+||r||\textasciicircum{}2<=C||r||\textasciicircum{}2 (GT-Pprops ||P||<=1+C eta; GT-bh-cstar ||r\textasciicircum{}2||=||r||\textasciicircum{}2). STEP B (shift to a positive element): set gamma=C delta||r||\textasciicircum{}2 the constant from node 1.2, so q\_r >= -gamma 1 (node 1.2), hence y:=q\_r+gamma 1 >= 0 (GT-orderunit-def). Also ||y||<=||q\_r||+gamma<=C||r||\textasciicircum{}2+C delta||r||\textasciicircum{}2<=C||r||\textasciicircum{}2 (Step A; gamma=O(delta||r||\textasciicircum{}2)). Thus 0 <= y <= ||y|| 1 <= M 1 with M:=C||r||\textasciicircum{}2 (GT-jb-orderunit-3310 / GT-orderunit-def: ||y||<=M and y>=0 give 0<=y<=M1). STEP C (Phi(y) is small and positive): P(q\_r)=0 (node 1.1), so Phi(q\_r)=(Phi-P)(q\_r), ||Phi(q\_r)||<=delta||q\_r||<=C delta||r||\textasciicircum{}2 (Step A). Phi(gamma 1)=gamma 1 (Phi unital), so Phi(y)=Phi(q\_r)+gamma 1 with ||Phi(y)||<=||Phi(q\_r)||+gamma<=C delta||r||\textasciicircum{}2. Also y>=0 => Phi(y)>=0 (GT-positive-monotone). STEP D (key: y\textasciicircum{}2<=M y): by node 1.3.1, 0<=y<=M1 gives 0<=y\textasciicircum{}2<=M y. STEP E (apply Phi, monotone): from Step D, 0<=y\textasciicircum{}2 and M y - y\textasciicircum{}2>=0; by GT-positive-monotone Phi(M y - y\textasciicircum{}2)>=0 i.e. Phi(y\textasciicircum{}2)<=M Phi(y), and Phi(y\textasciicircum{}2)>=0. Hence 0<=Phi(y\textasciicircum{}2)<=M Phi(y), so by norm-monotonicity (GT-sup-states: 0<=u<=v => ||u||<=||v|| in the order unit space B(H)\_sa) ||Phi(y\textasciicircum{}2)||<=||M Phi(y)||=M||Phi(y)||<=C||r||\textasciicircum{}2 * C delta||r||\textasciicircum{}2=C delta||r||\textasciicircum{}4. STEP F (back to q\_r\textasciicircum{}2): q\_r=y-gamma 1, so q\_r\textasciicircum{}2=(y-gamma1) o (y-gamma1)=y\textasciicircum{}2-2gamma y+gamma\textasciicircum{}2 1, and the Jordan identity gives 2y\textasciicircum{}2+2gamma\textasciicircum{}2 1 - q\_r\textasciicircum{}2=(y+gamma1)\textasciicircum{}2>=0 (a square is positive, GT-jb-positive-spectrum), i.e. q\_r\textasciicircum{}2<=2y\textasciicircum{}2+2gamma\textasciicircum{}2 1. Apply Phi (GT-positive-monotone) and triangle/norm-monotonicity: ||Phi(q\_r\textasciicircum{}2)||<=||Phi(2y\textasciicircum{}2+2gamma\textasciicircum{}2 1)||<=2||Phi(y\textasciicircum{}2)||+2gamma\textasciicircum{}2||Phi(1)||=2||Phi(y\textasciicircum{}2)||+2gamma\textasciicircum{}2 (Phi(1)=1). With ||Phi(y\textasciicircum{}2)||<=C delta||r||\textasciicircum{}4 (Step E) and gamma\textasciicircum{}2=C delta\textasciicircum{}2||r||\textasciicircum{}4=O(delta\textasciicircum{}2||r||\textasciicircum{}4)<=C delta||r||\textasciicircum{}4 (delta<=delta\_0=C eta\_0, eta\_0<1/4 fixed): ||Phi(q\_r\textasciicircum{}2)||<=C delta||r||\textasciicircum{}4. STEP G (replace Phi by P): ||(P-Phi)(q\_r\textasciicircum{}2)||<=delta||q\_r\textasciicircum{}2||=delta||q\_r||\textasciicircum{}2<=delta(C||r||\textasciicircum{}2)\textasciicircum{}2=C delta||r||\textasciicircum{}4 (Step A; GT-bh-cstar ||q\_r\textasciicircum{}2||=||q\_r||\textasciicircum{}2 for q\_r=q\_r\textasciicircum{}*). Hence ||P(q\_r\textasciicircum{}2)||<=||Phi(q\_r\textasciicircum{}2)||+C delta||r||\textasciicircum{}4<=C delta||r||\textasciicircum{}4<=C eta||r||\textasciicircum{}4 (delta<=C eta). Uses node 1.1, node 1.2, node 1.3.1. \steptag{validated} \\[6pt]
\step{1.3.1} FUNCTIONAL CALCULUS: for self-adjoint y in V=B(H)\_sa with 0<=y<=M 1 (M>=0), one has 0<=y\textasciicircum{}2<=M y. PROOF via the JB spectral theorem. V=B(H)\_sa is a unital JB algebra (GT-bhsa-jc + GT-jc-is-jb). By GT-jb-spectral (HOS 3.2.4), C(y) (smallest norm-closed Jordan subalgebra of V containing y and 1) is isometrically and Jordan-isomorphically mapped by phi onto C(Sp y), the continuous real functions on the compact spectrum Sp y subset R, with phi(1)=1 (constant function 1) and phi(y)=iota (the identity function t $|-\rangle$ t). (i) SPECTRUM IN [0,M]: y, M1-y, and y\textasciicircum{}2, M y - y\textasciicircum{}2 all lie in C(y) (Jordan subalgebra generated by y,1; y\textasciicircum{}2=y o y, M y=scalar mult, M1=scalar mult of 1). By GT-jb-positive-spectrum positivity is intrinsic (c>=0 iff Sp c subset [0,infty), spectrum computed in C(c) subset C(y)), so the hypotheses y>=0 and M1-y>=0 (from 0<=y<=M1 in V; M1-y>=0 is M1-y in V\_+, and M1-y in C(y)) transfer to C(y): under phi, phi(y)=t>=0 and phi(M1-y)=M-t>=0 pointwise on Sp y, hence Sp y subset [0,M]. (ii) POINTWISE INEQUALITIES: phi is a unital Jordan isomorphism, so phi(y\textasciicircum{}2)=phi(y)\textasciicircum{}2=t\textasciicircum{}2 and phi(M y)=M*phi(y)=M t (real-linear, multiplicative on the associative C(y)). On Sp y subset [0,M]: t\textasciicircum{}2>=0 and M t - t\textasciicircum{}2 = t(M-t)>=0 pointwise (product of two nonnegatives, t>=0 and M-t>=0). So phi(y\textasciicircum{}2)>=0 and phi(M y - y\textasciicircum{}2)=M t - t\textasciicircum{}2>=0 pointwise in C(Sp y). (iii) PULL BACK: the cone of C(Sp y) is the pointwise-nonnegative functions (an order unit space, GT-jb-positive-spectrum/HOS:474), and phi is an order isomorphism (positivity intrinsic, preserved by the Jordan isomorphism), so y\textasciicircum{}2=phi\textasciicircum{}\{-1\}(t\textasciicircum{}2)>=0 and M y - y\textasciicircum{}2=phi\textasciicircum{}\{-1\}(M t - t\textasciicircum{}2)>=0 in C(y), hence in V (positivity intrinsic, GT-jb-positive-spectrum). Therefore 0<=y\textasciicircum{}2<=M y in V. Uses GT-jb-spectral, GT-jb-positive-spectrum, GT-bhsa-jc, GT-jc-is-jb. \steptag{validated} \\[4pt]
\step{1.4} SEMINORM BOUND: ||q\_r||\_omega <= C sqrt(eta) ||r||\textasciicircum{}2 for every state seminorm omega=rho o Phi. By def-state-seminorm, ||q\_r||\_omega\textasciicircum{}2=omega(q\_r\textasciicircum{}2)=rho(Phi(q\_r\textasciicircum{}2)). rho is a state on B(H), so its restriction to B(H)\_sa is a state on the order-unit space B(H)\_sa (order unit 1), hence |rho(a)|<=||a|| for self-adjoint a (GT-state-bounded). Phi(q\_r\textasciicircum{}2) is self-adjoint (Phi maps B(H)\_sa to B(H)\_sa, q\_r\textasciicircum{}2 in B(H)\_sa). Therefore ||q\_r||\_omega\textasciicircum{}2=rho(Phi(q\_r\textasciicircum{}2))<=|rho(Phi(q\_r\textasciicircum{}2))|<=||Phi(q\_r\textasciicircum{}2)||. By node 1.3 (Step F), ||Phi(q\_r\textasciicircum{}2)||<=C delta||r||\textasciicircum{}4<=C eta||r||\textasciicircum{}4 (delta<=C eta, GT-Pprops). Hence ||q\_r||\_omega\textasciicircum{}2<=C eta||r||\textasciicircum{}4, and taking square roots ||q\_r||\_omega<=sqrt(C eta)||r||\textasciicircum{}2=C' sqrt(eta)||r||\textasciicircum{}2 with C'=sqrt(C) a universal constant. Uses node 1.3. \steptag{validated} \\[4pt]
\end{proofbox}

\end{document}
